% GENERATED FILE. Edit canonical lessons, then run npm run generate:books.
% canonical-source-hash: fnv1a64:879dd57c94a45709
% canonical-lessons: GU-C419-seri, GU-C419-phaliyu, GU-C419-mahollo, GU-C419-rahevasi, GU-C419-rajhalvu

\chapter{Streets and Lanes}
\label{ch:gu-streets-and-lanes}
\hlchaptermodality{419}

\begin{chapteropening}
\textbf{By the end of this chapter:} \emph{I can talk about my street and the people who live there.}

Complete the last lesson of chapter 419: I can talk about my street and the people who live there.
\end{chapteropening}

\section[\texorpdfstring{\gu{શેરી} (śerī)}{શેરી (śerī)}]{\gu{શેરી} (śerī) — a street; a lane (in a town)}

\label{lesson:GU-C419-seri}



\begin{quote}
Before the new one: say the Gujarati for a water tank, then the Gujarati for a tenant.
\end{quote}

\subsection*{You'll want to know: \gu{શેરી}}
\textbf{\gu{શેરી}} — \emph{śerī} — \textquotedblleft{}a street; a lane (in a town)\textquotedblright{}.

\begin{quote}
\gu{અમારી} \gu{શેરીમાં} \gu{બધાં} \gu{ઘર} \gu{જૂનાં} \gu{છે}.
\end{quote}

\begin{quote}
\emph{amārī śerīmā̃ badhā̃ ghar jūnā̃ che.}
\end{quote}

\begin{quote}
All the houses in our street are old.
\end{quote}

A feminine noun.

\begin{cousinweb}
\textbf{\gu{શેરી}} (\emph{śerī}) No dictionary we use records an older form, so its history is not settled.
\end{cousinweb}

\begin{grammarlens}[title={What you've built}]
One more word for getting around.
\end{grammarlens}

\subsection*{Your turn}
\begin{itemize}
\raggedright
  \item \emph{Say it:} \emph{śerī}
  \item \emph{Say it:} \emph{śerī}, once more
  \item \emph{Say it:} \emph{amārī śerīmā̃ badhā̃ ghar jūnā̃ che.} — All the houses in our street are old.
  \item \emph{Recall:} say the Gujarati for a water tank, then the Gujarati for a tenant, then say \emph{śerī} again
\end{itemize}

\begin{tcolorbox}[breakable,colback=teal!4,colframe=teal!35!black,title={Before you move on}]
What does \gu{શેરી} mean? (\textquotedblleft{}A street; a lane (in a town)\textquotedblright{}.) Say it once more.
\end{tcolorbox}
\section[\texorpdfstring{\gu{ફળિયું} (phaḷiyũ)}{ફળિયું (phaḷiyũ)}]{\gu{ફળિયું} (phaḷiyũ) — a neighbourhood lane (and its houses)}

\label{lesson:GU-C419-phaliyu}



\begin{quote}
Before the new one: say the Gujarati for a tenant, then the Gujarati for a street.
\end{quote}

\subsection*{You'll want to know: \gu{ફળિયું}}
\textbf{\gu{ફળિયું}} — \emph{phaḷiyũ} — \textquotedblleft{}a neighbourhood lane (and its houses)\textquotedblright{}.

\begin{quote}
\gu{ફળિયાનાં} \gu{બાળકો} \gu{સાથે} \gu{રમે} \gu{છે}.
\end{quote}

\begin{quote}
\emph{phaḷiyānā̃ bāḷko sāthe rame che.}
\end{quote}

\begin{quote}
The children of the lane play together.
\end{quote}

A neuter noun.

\begin{cousinweb}
\textbf{\gu{ફળિયું}} (\emph{phaḷiyũ}) is a lane together with the houses that open onto it, and the families in them. No dictionary we use records an older form, so its history is not settled.
\end{cousinweb}

\begin{grammarlens}[title={What you've built}]
One more word for getting around.
\end{grammarlens}

\subsection*{Your turn}
\begin{itemize}
\raggedright
  \item \emph{Say it:} \emph{phaḷiyũ}
  \item \emph{Say it:} \emph{phaḷiyũ}, once more
  \item \emph{Say it:} \emph{phaḷiyānā̃ bāḷko sāthe rame che.} — The children of the lane play together.
  \item \emph{Recall:} say the Gujarati for a tenant, then the Gujarati for a street, then say \emph{phaḷiyũ} again
\end{itemize}

\begin{tcolorbox}[breakable,colback=teal!4,colframe=teal!35!black,title={Before you move on}]
What does \gu{ફળિયું} mean? (\textquotedblleft{}A neighbourhood lane (and its houses)\textquotedblright{}.) Say it once more.
\end{tcolorbox}
\section[\texorpdfstring{\gu{મહોલ્લો} (mahollo)}{મહોલ્લો (mahollo)}]{\gu{મહોલ્લો} (mahollo) — a quarter of a town; a locality}

\label{lesson:GU-C419-mahollo}



\begin{quote}
Before the new one: say the Gujarati for a street, then the Gujarati for a neighbourhood lane.
\end{quote}

\subsection*{You'll want to know: \gu{મહોલ્લો}}
\textbf{\gu{મહોલ્લો}} — \emph{mahollo} — \textquotedblleft{}a quarter of a town; a locality\textquotedblright{}.

\begin{quote}
\gu{આ} \gu{મહોલ્લામાં} \gu{એક} \gu{મોટું} \gu{મંદિર} \gu{છે}.
\end{quote}

\begin{quote}
\emph{ā mahollāmā̃ ek moṭũ mandir che.}
\end{quote}

\begin{quote}
There is a big temple in this quarter.
\end{quote}

A masculine noun.

\begin{cousinweb}
\textbf{\gu{મહોલ્લો}} (\emph{mahollo}) is usually traced to Arabic \emph{maḥalla}, \textquotedblleft{}a place where people settle; a quarter of a town\textquotedblright{}, which came through Persian.
\end{cousinweb}

\begin{grammarlens}[title={What you've built}]
One more word for getting around.
\end{grammarlens}

\subsection*{Your turn}
\begin{itemize}
\raggedright
  \item \emph{Say it:} \emph{mahollo}
  \item \emph{Say it:} \emph{mahollo}, once more
  \item \emph{Say it:} \emph{ā mahollāmā̃ ek moṭũ mandir che.} — There is a big temple in this quarter.
  \item \emph{Recall:} say the Gujarati for a street, then the Gujarati for a neighbourhood lane, then say \emph{mahollo} again
\end{itemize}

\begin{tcolorbox}[breakable,colback=teal!4,colframe=teal!35!black,title={Before you move on}]
What does \gu{મહોલ્લો} mean? (\textquotedblleft{}A quarter of a town; a locality\textquotedblright{}.) Say it once more.
\end{tcolorbox}
\section[\texorpdfstring{\gu{રહેવાસી} (rahevāsī)}{રહેવાસી (rahevāsī)}]{\gu{રહેવાસી} (rahevāsī) — a resident; an inhabitant}

\label{lesson:GU-C419-rahevasi}



\begin{quote}
Before the new one: say the Gujarati for a neighbourhood lane, then the Gujarati for a quarter of a town.
\end{quote}

\subsection*{You'll want to know: \gu{રહેવાસી}}
\textbf{\gu{રહેવાસી}} — \emph{rahevāsī} — \textquotedblleft{}a resident; an inhabitant\textquotedblright{}.

\begin{quote}
\gu{ગામના} \gu{રહેવાસીઓએ} \gu{નવો} \gu{કૂવો} \gu{ખોદ્યો}.
\end{quote}

\begin{quote}
\emph{gāmnā rahevāsīoe navo kūvo khodyo.}
\end{quote}

\begin{quote}
The villagers dug a new well.
\end{quote}

A masculine or feminine noun, by the person meant.

\begin{cousinweb}
\textbf{\gu{રહેવાસી}} (\emph{rahevāsī}) is built from \gu{રહેવું}, to live, and -\gu{વાસી}, \textquotedblleft{}a dweller\textquotedblright{}, an ending that goes back to Sanskrit \emph{vāsa}, \textquotedblleft{}a dwelling\textquotedblright{} (Monier-Williams): the people who live in a place.
\end{cousinweb}

\begin{grammarlens}[title={What you've built}]
One more word for getting around.
\end{grammarlens}

\subsection*{Your turn}
\begin{itemize}
\raggedright
  \item \emph{Say it:} \emph{rahevāsī}
  \item \emph{Say it:} \emph{rahevāsī}, once more
  \item \emph{Say it:} \emph{gāmnā rahevāsīoe navo kūvo khodyo.} — The villagers dug a new well.
  \item \emph{Recall:} say the Gujarati for a neighbourhood lane, then the Gujarati for a quarter of a town, then say \emph{rahevāsī} again
\end{itemize}

\begin{tcolorbox}[breakable,colback=teal!4,colframe=teal!35!black,title={Before you move on}]
What does \gu{રહેવાસી} mean? (\textquotedblleft{}A resident; an inhabitant\textquotedblright{}.) Say it once more.
\end{tcolorbox}
\section[\texorpdfstring{\gu{રઝળવું} (rajhaḷvũ)}{રઝળવું (rajhaḷvũ)}]{\gu{રઝળવું} (rajhaḷvũ) — to wander aimlessly; to be left roaming}

\label{lesson:GU-C419-rajhalvu}



\begin{quote}
Before the new one: say the Gujarati for a quarter of a town, then the Gujarati for a resident.
\end{quote}

\subsection*{You'll want to know: \gu{રઝળવું}}
\textbf{\gu{રઝળવું}} — \emph{rajhaḷvũ} — \textquotedblleft{}to wander aimlessly; to be left roaming\textquotedblright{}.

\begin{quote}
\gu{રાત્રે} \gu{શેરીમાં} \gu{કૂતરા} \gu{રઝળે} \gu{છે}.
\end{quote}

\begin{quote}
\emph{rātre śerīmā̃ kūtrā rajhḷe che.}
\end{quote}

\begin{quote}
At night dogs wander the street.
\end{quote}

Past (he …): \gu{તે} \gu{આખો} \gu{દિવસ} \gu{રઝળ્યો}, \emph{te ākho divas rajhaḷyo}.

\begin{cousinweb}
\textbf{\gu{રઝળવું}} (\emph{rajhaḷvũ}) No dictionary we use records an older form, so its history is not settled.
\end{cousinweb}

\begin{grammarlens}[title={What you've built}]
One more word for getting around.
\end{grammarlens}

\subsection*{Your turn}
\begin{itemize}
\raggedright
  \item \emph{Say it:} \emph{rajhaḷvũ}
  \item \emph{Say it:} \emph{rajhaḷvũ}, once more
  \item \emph{Say it:} \emph{rātre śerīmā̃ kūtrā rajhḷe che.} — At night dogs wander the street.
  \item \emph{Recall:} say the Gujarati for a quarter of a town, then the Gujarati for a resident, then say \emph{rajhaḷvũ} again
\end{itemize}

\begin{tcolorbox}[breakable,colback=teal!4,colframe=teal!35!black,title={Before you move on}]
What does \gu{રઝળવું} mean? (\textquotedblleft{}To wander aimlessly; to be left roaming\textquotedblright{}.) Say it once more.
\end{tcolorbox}
